\documentclass[11pt]{article}
\usepackage[margin=1in]{geometry}
\usepackage[T1]{fontenc}
\usepackage{hyperref}
\usepackage{enumitem}
\usepackage{xcolor}
\usepackage{parskip}

\title{Marswalk: Project Plan and Handoff}
\author{Hackathon team}
\date{Deadline: 13 November 2026, midnight (local time)}

\begin{document}
\maketitle

\section{Goal}
Build an interactive map that a Mars astronaut could use to plan and carry out a surface excursion at Gale Crater. The map layers NASA data (optical imagery, topography, and steep-terrain hazard areas), shows science stops and a route, and gives crew-facing context for each stop.

\section{Current status}
\begin{itemize}[leftmargin=*]
  \item[\textcolor{green!50!black}{[done]}] Leaflet map in the Mars Trek equirectangular projection (EPSG:4326), centered on Gale Crater.
  \item[\textcolor{green!50!black}{[done]}] Base layers: Viking color mosaic (optical) and MOLA shaded relief (topography).
  \item[\textcolor{green!50!black}{[done]}] Base point and three science stops, each with a summary and objective.
  \item[\textcolor{green!50!black}{[done]}] Dashed route from base through the stops and back.
  \item[\textcolor{green!50!black}{[done]}] Steep-terrain hazard layer computed from MOLA elevation (slope over $15^\circ$), saved as \texttt{steep.geojson}.
  \item[\textcolor{orange!80!black}{[todo]}] Confirm \texttt{steep.geojson} loads and the hazard areas appear on the map.
  \item[\textcolor{orange!80!black}{[todo]}] Route does not yet avoid hazard areas.
  \item[\textcolor{orange!80!black}{[todo]}] No Curiosity measurements (chemistry, radiation) attached to stops yet.
\end{itemize}

\section{Repository layout}
\begin{description}[leftmargin=2.5cm, style=nextline]
  \item[\texttt{src/App.tsx}] Main page: map, layer controls, stop sidebar, route, and hazard layer.
  \item[\texttt{src/data/site.ts}] Mission data: base point, stops (coordinates and text), route waypoints.
  \item[\texttt{src/App.css}, \texttt{src/index.css}] Styling. \texttt{index.css} must contain only the short global block (no template rules).
  \item[\texttt{public/steep.geojson}] Hazard polygons. Must sit directly in \texttt{public/}, not a subfolder.
  \item[\texttt{docs/PROJECT\_PLAN.tex}] This document.
\end{description}

\section{Running the app}
\begin{enumerate}[leftmargin=*]
  \item Install Node.js (LTS version) if it is not installed.
  \item Open a terminal in the project folder and run:
\begin{verbatim}
npm install
npm run dev
\end{verbatim}
  \item Open the local address printed in the terminal (usually \texttt{http://localhost:5173}).
  \item Check that \texttt{http://localhost:5173/steep.geojson} shows JSON text. If it shows ``Not Found,'' the file is in the wrong place.
\end{enumerate}

\section{Hazard data pipeline (already run once)}
The steep-terrain file was produced with GDAL (we used the conda-forge build on Windows). Re-run these steps only if you change the crop area or the threshold.

\begin{enumerate}[leftmargin=*]
  \item Download the MOLA elevation mosaic (about 2 GB):
\begin{verbatim}
curl -L -O https://planetarymaps.usgs.gov/mosaic/Mars_MGS_MOLA_DEM_mosaic_global_463m.tif
\end{verbatim}
  \item Crop to Gale Crater. The file is in meters, not degrees. The box is 134--141$^\circ$E, 2--8$^\circ$S, using about 59{,}276 m per degree:
\begin{verbatim}
gdal_translate -projwin 7943051 -118553 8358986 -474212 Mars_MGS_MOLA_DEM_mosaic_global_463m.tif gale.tif
\end{verbatim}
  \item Compute slope in degrees (no \texttt{-s} option, because the file is already in meters):
\begin{verbatim}
gdaldem slope gale.tif slope.tif -compute_edges
\end{verbatim}
  \item Mark cells steeper than $15^\circ$:
\begin{verbatim}
python "C:\Users\HP\miniconda3\envs\gdal\Scripts\gdal_calc.py" -A slope.tif --calc="A>15" --outfile=steep.tif --type=Byte --NoDataValue=0
\end{verbatim}
  \item Convert marked cells to polygons:
\begin{verbatim}
gdal raster polygonize steep.tif steep.geojson
\end{verbatim}
  \item Copy \texttt{steep.geojson} into \texttt{public/}.
\end{enumerate}

\textbf{Limit to state on the map:} at 463 m per pixel, the layer shows broad steep terrain, not individual rocks. Individual hazards would need high-resolution imagery such as HiRISE.

\section{Known issues}
\begin{itemize}[leftmargin=*]
  \item \textbf{Viking layer URL.} The Viking tile ID was taken from the Mars Trek API page. It must be checked in a browser. Test: \texttt{https://trek.nasa.gov/tiles/Mars/EQ/Mars\_Viking\_MDIM21\_ClrMosaic\_global\_232m/1.0.0/WMTSCapabilities.xml}.
  \item \textbf{Stop coordinates are estimates.} Check them against NASA's published Curiosity locations before submitting.
  \item \textbf{Hazard fetch fails silently.} If \texttt{steep.geojson} is missing, the map shows no hazards and no error. Check the URL in the first place.
  \item \textbf{Remote access.} The Mars Trek tiles load from the user's browser, not from the build. Some networks may block them.
\end{itemize}

\section{Remaining tasks, in priority order}
\begin{enumerate}[leftmargin=*]
  \item Confirm the hazard layer displays (see Known issues). Zoom to Base to check.
  \item Confirm both base layers load.
  \item Add real science content to each stop: Curiosity measurements or published results (for example chemistry or radiation), with sources cited on the page.
  \item Make the route avoid hazard areas, or add a warning when a leg crosses one.
  \item Add a crew-facing note for each stop (what the crew will see, what to do there).
  \item Prepare the presentation: state the data sources, the $15^\circ$ threshold, and the 463 m resolution limit.
  \item Commit everything, including \texttt{public/steep.geojson}, and push. Note that GitHub may need the Claude GitHub App installed for the owner account before pushing from this environment.
\end{enumerate}

\section{Notes for the next person}
\begin{itemize}[leftmargin=*]
  \item Most of the logic lives in \texttt{src/App.tsx} and the stop data lives in \texttt{src/data/site.ts}. Start there.
  \item The map uses an equirectangular projection (\texttt{L.CRS.EPSG4326}). Do not switch it to the default web-mercator projection, or the Mars tiles will not line up.
  \item Coordinates in \texttt{site.ts} are \texttt{[latitude, longitude]}. Leaflet uses the same order.
\end{itemize}

\end{document}
